% Book Two Pass 2: discrete three-stage geodesics-in-heat reproduction.
\clearpage
\section*{Research reconstruction: heat flow becomes surface distance}
\noindent\textbf{Primary work:} Keenan Crane, Clarisse Weischedel and
Max Wardetzky,
\href{https://www.cs.cmu.edu/~kmcrane/Projects/HeatMethod/}
{\emph{Geodesics in Heat} (2013)}. The idea is surprising:
a short-time diffusion computation can reveal the direction
toward a geodesic distance field, although heat itself is not
geodesic distance.

\paragraph{Stage I: a source becomes a temperature field.}
Let \(K\) be the positive-semidefinite finite-element stiffness
and \(M\) the mass matrix. With our sign convention, a backward
Euler step of the heat equation is
\[
(M+tK)u=M u_0,\qquad t>0.
\]
The initial vector \(u_0\) represents a source impulse in the
chosen discrete model. The positive timestep stabilizes the
linear solve. Choosing a timestep tied to squared mean edge
length is a practical resolution-dependent heuristic, not a
universal accuracy proof.

\paragraph{Stage II: temperature supplies a direction.}
On each triangle, compute the piecewise-constant gradient of
\(u\) from its three nodal values and element basis gradients.
Then define the unit vector field
\[
X=-\frac{\nabla u}{\lVert\nabla u\rVert}.
\]
The negative sign makes \(X\) point outward, away from a hot
source. If \(\lVert\nabla u\rVert\) is too small to normalize
reliably, the direction is undefined and the code must handle
that case, rather than divide by a near-zero value.

\paragraph{Stage III: integrate those directions.}
Find a scalar field whose gradient approximates \(X\) by
minimizing \(\frac12\int\lVert\nabla\phi-X\rVert^2dA\).
The first-order condition is the Poisson system
\[
K\phi=b,\qquad
b_i=\sum_{T\ni i}A_T\,\nabla\phi_i\cdot X_T.
\]
Because a constant can be added to \(\phi\) without changing
its gradient, \(K\) has a constant-vector nullspace on a
free-boundary connected mesh. Fixing \(\phi\) to zero at the
source selects a coordinate gauge; it does not create a new
distance observation.

\paragraph{Independent analytical oracle.}
On a flat rectangular convex domain with a central source and
no obstacles, the shortest intrinsic path to another interior
point is the straight line. Its Euclidean length
\(\lVert p-p_{\mathrm{source}}\rVert\) therefore supplies a reference
for a simple grid triangulation. The companion code constructs
its own P1 stiffness/mass matrices, takes one backward heat
step, normalizes triangle gradients, solves the anchored Poisson
equation, and compares its distances with this exact planar
reference.

\paragraph{What the method does \emph{not} claim.}
This teaching experiment does not reproduce the full paper's
large-mesh benchmark figures, curved-surface comparisons,
intrinsic-Delaunay corrections or production factorization
performance. It also does not establish accurate geodesics on
nonmanifold input, disconnected domains, degenerate triangles
or arbitrary boundary models. Those require additional algorithms
and domain checks.

\paragraph{Reconstruction and research questions.}
First demonstrate mass conservation after the implicit heat
step, stiffness symmetry, constant-mode behavior and the
Poisson gauge. Next double the grid resolution, vary \(t/h^2\),
and report median and maximum \emph{absolute} geodesic error
separately. Finally compare the heat method with a graph
shortest-path algorithm: graph edge distances are a useful
baseline but are not a ground truth for smooth surface distance.
A reliable figure must show both accuracy and mesh resolution,
not merely a colorful geodesic gradient.
